\section{Profesní profil}
Softwarový vývojář s více než třemi lety praxe v interních podnikových aplikacích a integracích. Propojuji analýzu požadavků s návrhem, full-stack vývojem a nasazením řešení v C\#/.NET a TypeScriptu. Spolupracuji s mezinárodními týmy a koordinuji dodávky softwaru od externích partnerů v týdenních sprintech. Mezi mé znalosti patří Codex, ChatGPT, LLM API a MCP. Ve společnosti 21 Consult Group chci propojit vývojářskou praxi se zájmem o AI při tvorbě interních nástrojů, automatizací a nových produktů.

\section{Pracovní zkušenosti}
\cvrole{KAMAX}{Software Engineer}{06/2022–08/2024; 01/2025–dosud}
\begin{highlights}
  \item Analýza požadavků, návrh řešení a full-stack vývoj interních aplikací v C\#/.NET, TypeScriptu, Reactu a Angularu s databází Microsoft SQL Server.
  \item Nástroje pro řízení kvality a výrobní logistiku včetně aplikací pro obsluhu vysokozdvižných vozíků.
  \item Koordinace externího dodavatele při zavádění MES: týdenní sprinty, kontrola postupu a řešení otevřených bodů.
  \item Integrace s podnikovými a výrobními systémy MES, SAP a SharePoint.
  \item Správa GitLabu, repozitářů a CI/CD pipeline; nasazování aplikací prostřednictvím IIS a Dockeru.
  \item Mezinárodní projekty a spolupráce s německými závody; pracovní komunikace v angličtině a v některých případech vedení projektu.
\end{highlights}

\vspace{5pt}
\cvrole{ioLabs}{Software Developer\enspace/\enspace externí kontrakt}{09/2024–11/2024}
\begin{highlights}
  \item Vývoj a údržba aplikací v ASP.NET/.NET s databází Microsoft SQL Server v rámci externího kontraktu.
\end{highlights}

\section{Vzdělání}
\cvrole{Technická univerzita v Liberci}{Aplikovaná informatika\enspace/\enspace Ing., s vyznamenáním}{2022–2024}
{\small Diplomová práce: Systém pro predikci a vytěžování výroby FVE.}\par
\vspace{7pt}
\cvrole{Technická univerzita v Liberci}{Aplikovaná informatika\enspace/\enspace bakalářské studium}{2019–2022}
{\small Bakalářská práce: Identifikace indických jazyků s využitím neuronových sítí.}\par

\newpage
\section{Technické znalosti}
\skillrow{AI nástroje}{%
  Codex, ChatGPT.
}
\skillrow{AI a LLM}{%
  LLM API, prompt engineering, MCP servery, LangChain, LangGraph, smolagents; RAG (znalost principů). Obecné principy návrhu AI agentů a propojování LLM s aplikacemi a nástroji.
}
\skillrow{Backend a API}{%
  C\#, ASP.NET Core, Entity Framework Core, návrh REST API, OpenAPI/Swagger.
}
\skillrow{Frontend}{%
  JavaScript, TypeScript, React, Angular, HTML5/CSS3.
}
\skillrow{Databáze}{%
  Microsoft SQL Server, PostgreSQL, MongoDB.
}
\skillrow{Git a provoz}{%
  Git/GitLab, GitLab CI/CD, Docker, správa Linux serverů, IIS/Nginx.
}
\skillrow{Bezpečnost}{%
  OAuth 2.0/OpenID Connect, Keycloak/OpenIddict, SSO.
}
\skillrow{Architektura a testy}{%
  Clean Architecture/DDD, unit testy (xUnit/NUnit).
}
\skillrow{Python}{%
  Školní projekty a AI projekty v rámci kurzů.
}

\section{Jazyky}
\skillrow{Čeština}{Rodilý mluvčí.}
\skillrow{Angličtina}{B2; pracovní komunikace v mezinárodních týmech.}
